\vspace{6pt}
\begin{tcolorbox}[abstractstyle,
  title={执行摘要 \quad {\normalsize\normalfont\itshape (Executive Summary)}}]

纠错码（Error-Correcting Code, ECC）是贯穿现代信息社会全域的核心基础技术——无论是5G/6G无线通信、相干光纤传输、NAND Flash固态存储、量子计算系统，还是新型存算一体器件，无一例外地高度依赖纠错码保障信息可靠传输与处理。然而，尽管应用场景千差万别，当前ECC设计范式却被三大结构性困境所困：\textbf{（一）过度依赖专家经验}——从码型构造、度分布优化到译码算法选择，全程仰赖跨越信息论、信号处理与VLSI设计的深厚专家知识壁垒；\textbf{（二）优化严重依赖大规模仿真}——获取一条可靠的BER性能曲线须耗费数小时乃至数天的蒙特卡洛仿真，使大规模码型搜索在工程上几乎不可行；\textbf{（三）算法优化与硬件实现孤立割裂}——译码算法与芯片硬件设计两条并行轨道互不知情，导致算法确定后硬件端大量推倒重来，严重拖慢标准化进程。

\vspace{4pt}
本提案提出 \textbf{PrimeCoDec}——一个面向信息系统纠错码需求的\textbf{全链路AI自动化框架}（Full-Pipeline AI Co-Design for Coding and Decoding）。PrimeCoDec 将上述三大困境同时消解：用户仅需以自然语言输入原始需求，框架即可\textbf{全自动生成}从纠错码编码方案、译码算法、硬件部署方案直至行业标准草案的完整输出，全程无需专家干预。

\vspace{4pt}
框架的核心创新体现于四个均衡配置的技术层次：\textbf{（一）编码方案AI自动化}，以AI驱动的码型搜索替代专家主导的码结构设计，实现面向任意物理信道的最优码型自动生成；\textbf{（二）无仿真译码算法优化}，以DeepONet神经算子在极低延迟内预测完整BER/FER性能曲线，从根本上替代耗时的蒙特卡洛仿真，实现译码算法参数的快速自动调优；\textbf{（三）算法--硬件协同自动实现}，以联合代理模型在统一目标函数下同步优化算法性能与VLSI指标，支持共模架构（一套硬件实现多种码型译码）等高价值设计输出；\textbf{（四）顶层需求形式化与标准自动输出}，以大语言模型（LLM）将自然语言需求转化为形式化约束，并将优化结果逆向映射为符合3GPP/IEEE格式的行业标准草案。

\vspace{4pt}
作为核心验证场景，我们选择\textbf{5G NR标准（3GPP Release 17）}作为基准——这是当前工业界极少数经数十年专家深度打磨的完备行业标准之一，对比天然公平且极具说服力。PrimeCoDec将全自动生成Polar码与LDPC码的最优方案，给出当前场景下的最优译码算法，设计共模架构译码器（一套硬件同时实现Polar码与LDPC码译码），并输出修订后的3GPP编码标准草案——从用户输入到可部署标准，全程由框架自动完成。
\end{tcolorbox}
\vspace{4pt}